% Explanation certificate for heart-paper (genome 9022a9eb-ecf9-41ba-b138-28ed4d131c9d).
% Generated by explaneat export (schema 1.0.0) from the JSON record alone;
% content sha256 ae2a73baae0507a6bb122be13c78ec321d84882c02c55591f1d760ead16f486b.
% Drafter recorded in the JSON: unknown (None). Do not edit by hand: re-render with
%   uv run python -m explaneat export render --from heart-paper.json
% TikZ modules are input from \explaneatcertdir{<label>} (empty by default: same
% directory; index.tex sets it to the genome's subdirectory).
\providecommand{\explaneatcertdir}[1]{}
\begin{drafted}[appendix: heart-paper explanation certificate]
\section{Explanation certificate: Heart, AAAI paper genome}\label{app:heart-paper:sec}
This section is the explanation certificate of the Heart, AAAI paper genome network. It is rendered from the canonical JSON record \texttt{heart-paper.json} (schema 1.0.0, SHA-256 \texttt{ae2a73baae0507a6}\ldots), which also holds every evidence payload, the full graphs with weights, and the explanation record. \Cref{app:heart-paper:summary} summarises the network and its coverage.

\begin{table}[htbp]
\centering\small
\begin{tabularx}{\linewidth}{@{}lX@{}}
\toprule
Dataset & Heart Disease (binary, prepared) \\
Experiment & Heart Binary \\
Genome & \texttt{9022a9eb-ecf9-41ba-b138-28ed4d131c9d} (NEAT key 9493, generation 38) \\
Explanation & \texttt{a03ebc86-a14d-4860-9eb6-29f040d16633} (40 operations) \\
Train AUC / accuracy & 0.9455 / 0.8636 (n = 242) \\
Test AUC / accuracy & 0.8966 / 0.8033 (n = 61) \\
Output activation & sigmoid \\
Phenotype & 32 nodes (26 inputs, 5 hidden, 1 output), 31 enabled edges \\
Explained model state & 42 nodes, 34 enabled edges \\
Refactoring operations & split\_\allowbreak{}node 8 \\
Function preservation & max $|\Delta\,\mathrm{output}|$ = $2.22\times 10^{-16}$ on the training split \\
Annotations & 6 (4 leaf, 2 composition) \\
Structural coverage & 1 (41 of 41 candidate nodes; 7 Null) \\
Compositional coverage & 1 \\
Uncovered nodes / edges & 0 / 0 \\
Missing compositions & 0 \\
Full coverage & yes \\
Drafted by & unknown \\
Verified by & unverified \\
Record hash (SHA-256) & \texttt{ae2a73baae0507a6}\ldots \\
\bottomrule
\end{tabularx}
\caption{Explanation certificate summary: Heart, AAAI paper genome.}
\label{app:heart-paper:summary}
\end{table}

\subsection{Coverage}\label{app:heart-paper:coverage}
Coverage is computed by the current code (the function behind the \texttt{get\_coverage} tool) under these conventions. \emph{Candidate nodes} are all nodes of the explained model state except the output nodes. A candidate node is \emph{covered} when some annotation of the hierarchy (a leaf, or a composition together with its descendants) contains it and all its enabled outgoing edges. \emph{Exact Null} parts, which have no effect on any output for any input, count as covered: nodes with no path to an output over enabled, non-zero-weight edges (disconnected inputs and dead ends), and edges that have zero weight or touch such a node. Compositional coverage is the fraction of required compositions present: every annotation except a single global root containing every output must have a parent; a parentless, childless annotation made only of Null nodes is exempt.

\begin{itemize}
\item \textbf{Null, disconnected input (7):} \texttt{-12} (restecg:0), \texttt{-14} (restecg:2), \texttt{-17} (exang:1), \texttt{-21} (slope:3), \texttt{-23} (thal:3), \texttt{-3} (sex:1), \texttt{-6} (cp:3).
\item \textbf{Uncovered nodes:} none.
\item \textbf{Uncovered edges:} none.
\item \textbf{Missing compositions:} none.
\end{itemize}

\subsection{Annotation hierarchy}\label{app:heart-paper:hierarchy}
The hierarchy has 1 root; the global root is heart\_\allowbreak{}disease\_\allowbreak{}risk\_\allowbreak{}score. Tags (A1, A2, \ldots) are used in the figures.
\begingroup\parindent=0pt\parskip=1pt
\hspace*{0.0em}$\triangleright$~A1: heart\_\allowbreak{}disease\_\allowbreak{}risk\_\allowbreak{}score (composition, \cref{app:heart-paper:ann-heart-disease-risk-score})\par
\hspace*{1.5em}$\triangleright$~A2: benign\_\allowbreak{}profile\_\allowbreak{}gate (leaf, \cref{app:heart-paper:ann-benign-profile-gate})\par
\hspace*{1.5em}$\triangleright$~A3: restbp\_\allowbreak{}chol\_\allowbreak{}interaction (leaf, \cref{app:heart-paper:ann-restbp-chol-interaction})\par
\hspace*{1.5em}$\triangleright$~A4: risk\_\allowbreak{}aggregation\_\allowbreak{}gate (composition, \cref{app:heart-paper:ann-risk-aggregation-gate})\par
\hspace*{3.0em}$\triangleright$~A5: benign\_\allowbreak{}profile\_\allowbreak{}gate\_\allowbreak{}534path (leaf, \cref{app:heart-paper:ann-benign-profile-gate-534path})\par
\hspace*{3.0em}$\triangleright$~A6: chol\_\allowbreak{}threshold (leaf, \cref{app:heart-paper:ann-chol-threshold})\par
\endgroup

\subsection{Network figures}\label{app:heart-paper:figures}
\begin{figure}[htbp]
\centering
\input{\explaneatcertdir{heart-paper}heart-paper-fig-full.tex}
\caption{Heart, AAAI paper genome: full view of the explained model state.}
\label{app:heart-paper:fig-full}
\end{figure}
\begin{figure}[htbp]
\centering
\input{\explaneatcertdir{heart-paper}heart-paper-fig-level-0.tex}
\caption{Heart, AAAI paper genome: collapsed view at hierarchy level 0 (fully recomposed).}
\label{app:heart-paper:fig-level-0}
\end{figure}
\begin{figure}[htbp]
\centering
\input{\explaneatcertdir{heart-paper}heart-paper-fig-level-1.tex}
\caption{Heart, AAAI paper genome: collapsed view at hierarchy level 1.}
\label{app:heart-paper:fig-level-1}
\end{figure}
\begin{figure}[htbp]
\centering
\input{\explaneatcertdir{heart-paper}heart-paper-fig-level-2.tex}
\caption{Heart, AAAI paper genome: collapsed view at hierarchy level 2.}
\label{app:heart-paper:fig-level-2}
\end{figure}

\subsection{Refactoring operations}\label{app:heart-paper:operations}
The explained model state is the phenotype after these operations, replayed in order. All are function-preserving; the largest change in the model output over the 242 training rows is $2.22\times 10^{-16}$ (tolerance $1.00\times 10^{-9}$).
\begin{itemize}
\item \#26: split \texttt{-22} into \texttt{-22\_\allowbreak{}a}, \texttt{-22\_\allowbreak{}b}.
\item \#27: split \texttt{-9} into \texttt{-9\_\allowbreak{}a}, \texttt{-9\_\allowbreak{}b}, \texttt{-9\_\allowbreak{}c}, \texttt{-9\_\allowbreak{}d}.
\item \#28: split \texttt{-2} into \texttt{-2\_\allowbreak{}a}, \texttt{-2\_\allowbreak{}b}.
\item \#29: split \texttt{-10} into \texttt{-10\_\allowbreak{}a}, \texttt{-10\_\allowbreak{}b}.
\item \#30: split \texttt{155} into \texttt{155\_\allowbreak{}a}, \texttt{155\_\allowbreak{}b}.
\item \#31: split \texttt{-16} into \texttt{-16\_\allowbreak{}a}, \texttt{-16\_\allowbreak{}b}.
\item \#32: split \texttt{-2\_\allowbreak{}b} into \texttt{-2\_\allowbreak{}c}, \texttt{-2\_\allowbreak{}d}.
\item \#33: split \texttt{-10\_\allowbreak{}b} into \texttt{-10\_\allowbreak{}c}, \texttt{-10\_\allowbreak{}d}.
\end{itemize}

\subsection{Annotations}\label{app:heart-paper:annotations}
\subsubsection{\texorpdfstring{A1: heart\_\allowbreak{}disease\_\allowbreak{}risk\_\allowbreak{}score}{A1: heart disease risk score}}\label{app:heart-paper:ann-heart-disease-risk-score}
% annotation id ann_39, operation seq 39
\textbf{Kind:} Composition annotation (depth 0); children A2 (\cref{app:heart-paper:ann-benign-profile-gate}), A3 (\cref{app:heart-paper:ann-restbp-chol-interaction}), A4 (\cref{app:heart-paper:ann-risk-aggregation-gate}).

\begin{itemize}
\item \textbf{Entry nodes:} \texttt{-1} (age), \texttt{-4} (cp:1), \texttt{-5} (cp:2), \texttt{-7} (cp:4), \texttt{-11} (fbs:1), \texttt{-13} (restecg:1), \texttt{-15} (thalach), \texttt{-18} (oldpeak), \texttt{-19} (slope:1), \texttt{-20} (slope:2), \texttt{-24} (thal:4), \texttt{-25} (thal:6), \texttt{-26} (thal:7), \texttt{-22\_\allowbreak{}a} (ca\_\allowbreak{}a), \texttt{-9\_\allowbreak{}a} (chol\_\allowbreak{}a), \texttt{-2\_\allowbreak{}a} (sex:0\_\allowbreak{}a), \texttt{-10\_\allowbreak{}a} (fbs:0\_\allowbreak{}a), \texttt{-3} (sex:1), \texttt{-6} (cp:3), \texttt{-12} (restecg:0), \texttt{-14} (restecg:2), \texttt{-17} (exang:1), \texttt{-21} (slope:3), \texttt{-23} (thal:3).
\item \textbf{Exit nodes:} \texttt{0}.
\item \textbf{Internal nodes:} none.
\item \textbf{Edges (weight):} \texttt{-1}$\to$\texttt{0}~(0.0907), \texttt{-4}$\to$\texttt{0}~(-0.14), \texttt{-5}$\to$\texttt{0}~(0.683), \texttt{-7}$\to$\texttt{0}~(1.21), \texttt{-11}$\to$\texttt{0}~(0.0269), \texttt{-13}$\to$\texttt{0}~(0.19), \texttt{-15}$\to$\texttt{0}~(-0.256), \texttt{-18}$\to$\texttt{0}~(0.559), \texttt{-19}$\to$\texttt{0}~(-0.526), \texttt{-20}$\to$\texttt{0}~(0.853), \texttt{-24}$\to$\texttt{0}~(0.482), \texttt{-25}$\to$\texttt{0}~(-0.305), \texttt{-26}$\to$\texttt{0}~(0.463), \texttt{-22\_\allowbreak{}a}$\to$\texttt{0}~(2.03), \texttt{-9\_\allowbreak{}a}$\to$\texttt{0}~(0.676), \texttt{-2\_\allowbreak{}a}$\to$\texttt{0}~(-0.711), \texttt{-10\_\allowbreak{}a}$\to$\texttt{0}~(0.0357), \texttt{155\_\allowbreak{}a}$\to$\texttt{0}~(-0.96), \texttt{2379}$\to$\texttt{0}~(0.364), \texttt{534}$\to$\texttt{0}~(-0.884).
\item \textbf{With its descendants:} 42 nodes.
\item \textbf{Covers} (coverage, own and as part of the hierarchy): 41 candidate nodes.
\end{itemize}

\textbf{Hypothesis.} Root: y0 = ReLU(bias 0.318 + linear core + 3 nonlinear pathways). (1) DIRECT LINEAR CORE of 17 inputs to the output - ca (largest, +2.03), cp:4 (+1.21), slope:2, chol, cp:2, oldpeak push risk up; sex:0 (-0.71), slope:1, thalach push down. By SHAP the top drivers are ca, cp:4, sex:0, slope:2, chol. (2) benign\_\allowbreak{}profile\_\allowbreak{}gate (155\_\allowbreak{}a, -0.96) lowers risk for females without exertional angina. (3) restbp\_\allowbreak{}chol\_\allowbreak{}interaction (2379, +0.364). (4) risk\_\allowbreak{}aggregation\_\allowbreak{}gate (534, -0.884), the elaborate composed unit. 7 features unused. PERFORMANCE: train AUC 0.946 / test 0.897. ABLATION LADDER (test AUC, baseline 0.897): linear-core only 0.860 (nonlinearity buys +0.037); remove benign gate -0.028; remove restbp -0.012; remove 534 gate +0.002. So the two simple gates (benign, restbp) carry all the generalising benefit; the elaborate 534 aggregation gate is train-only. The benign signal is redundantly encoded twice (155\_\allowbreak{}a direct + 155\_\allowbreak{}b inside 534).

\textbf{Closed form} (children as named functions):
{\small\begingroup\allowdisplaybreaks
\begin{align*}
y_{0} &= \sigma\bigl(z_{0}\bigr) \\
z_{0} &= 0.318 + 0.676 x_{chol\_a} + 0.19 x_{restecg:1} + 0.482 x_{thal:4} + 2.03 x_{ca\_a} \\
&\qquad + 0.559 x_{oldpeak} + 0.0357 x_{fbs:0\_a} + 0.683 x_{cp:2} + 1.21 x_{cp:4} + 0.0907 x_{age} \\
&\qquad + 0.364 \operatorname{restbp\_chol\_interaction}\bigl(x_{trestbps}, x_{chol\_c}\bigr) \\
&\qquad + 0.853 x_{slope:2} + 0.0269 x_{fbs:1} + 0.463 x_{thal:7} - 0.256 x_{thalach} \\
&\qquad - 0.526 x_{slope:1} - 0.96 \operatorname{benign\_profile\_gate}\bigl(x_{sex:0\_b\_c}, \\
&\qquad x_{fbs:0\_b\_c}, x_{exang:0\_a}\bigr) - 0.14 x_{cp:1} - 0.305 x_{thal:6} - 0.711 x_{sex:0\_a} \\
&\qquad - 0.884 \operatorname{risk\_aggregation\_gate}\bigl(x_{ca\_b}, x_{chol\_b}, x_{sex:0\_b\_d}, \\
&\qquad x_{fbs:0\_b\_d}, x_{exang:0\_b}, x_{chol\_d}\bigr)
\end{align*}\endgroup}

Expanded over the children:
{\small\begingroup\allowdisplaybreaks
\begin{align*}
y_{0} &= \sigma\bigl(z_{0}\bigr) \\
z_{0} &= 0.318 + 0.676 x_{chol\_a} + 0.19 x_{restecg:1} + 0.482 x_{thal:4} + 2.03 x_{ca\_a} \\
&\qquad + 0.559 x_{oldpeak} + 0.0357 x_{fbs:0\_a} + 0.683 x_{cp:2} + 1.21 x_{cp:4} + 0.0907 x_{age} \\
&\qquad + 0.364 \operatorname{relu}\bigl(-0.754 + 0.931 \operatorname{relu}\bigl(3.4 \\
&\qquad + 0.978 x_{trestbps}\bigr) - 0.529 x_{chol\_c}\bigr) + 0.853 x_{slope:2} + 0.0269 x_{fbs:1} \\
&\qquad + 0.463 x_{thal:7} - 0.256 x_{thalach} - 0.526 x_{slope:1} \\
&\qquad - 0.96 \operatorname{relu}\bigl(0.0879 + 1.11 x_{exang:0\_a} + 0.824 x_{sex:0\_b\_c} \\
&\qquad - 0.509 x_{fbs:0\_b\_c}\bigr) - 0.14 x_{cp:1} - 0.305 x_{thal:6} - 0.711 x_{sex:0\_a} \\
&\qquad - 0.884 \operatorname{relu}\bigl(-0.761 + 0.818 \operatorname{relu}\bigl(0.0879 \\
&\qquad + 1.11 x_{exang:0\_b} + 0.824 x_{sex:0\_b\_d} - 0.509 x_{fbs:0\_b\_d}\bigr) \\
&\qquad - 0.592 x_{chol\_b} - 1.16 x_{ca\_b} - 0.677 \operatorname{relu}\bigl(-0.864 \\
&\qquad + 1.05 x_{chol\_d}\bigr)\bigr)
\end{align*}\endgroup}

where $y_{0}$ is node \texttt{0}; $z_{i}$, where shown, is the pre-activation of a sigmoid output $y_{i}$; weights and biases are shown to 3 significant figures.
 On the training split (n = 242), the displayed closed form differs from the forward pass by at most 0.00241, and the exact function by at most $2.22\times 10^{-16}$.

\textbf{Evidence} (2 items; full payloads in the JSON record):
\begin{itemize}
\item \textsc{empirical}, \texttt{performance\_\allowbreak{}metrics}; data split: train+test (inferred, inferred from payload keys). Whole-model performance and coverage. Complete structural and compositional coverage reached with 8 node splits and ZERO function-changing operations - the explained network is computationally identical to the trained one. Wider generalisation gap (0.049 AUC) on the small 61-row test split.
\item \textsc{empirical}, \texttt{ablation\_\allowbreak{}study}; data split: test (inferred, inferred from narrative). Full ablation ladder, re-verified end-to-end on the test split. Nonlinearity as a whole buys +0.037 test AUC, but the benefit comes almost entirely from the two shallow gates encoding recognised protective/risk factors. The elaborate composed gate contributes nothing on held-out data. A complete explanation doubles as a diagnosis of how much nonlinear structure earns its place.
\end{itemize}

\textbf{Authorship:} drafted by unknown on 2026-07-08; verified by nobody (unverified).

\subsubsection{\texorpdfstring{A2: benign\_\allowbreak{}profile\_\allowbreak{}gate}{A2: benign profile gate}}\label{app:heart-paper:ann-benign-profile-gate}
% annotation id ann_34, operation seq 34
\textbf{Kind:} Leaf annotation (depth 1); parent A1 (\cref{app:heart-paper:ann-heart-disease-risk-score}).

\begin{itemize}
\item \textbf{Entry nodes:} \texttt{-2\_\allowbreak{}c} (sex:0\_\allowbreak{}b\_\allowbreak{}c), \texttt{-10\_\allowbreak{}c} (fbs:0\_\allowbreak{}b\_\allowbreak{}c), \texttt{-16\_\allowbreak{}a} (exang:0\_\allowbreak{}a).
\item \textbf{Exit nodes:} \texttt{155\_\allowbreak{}a}.
\item \textbf{Internal nodes:} none.
\item \textbf{Edges (weight):} \texttt{-2\_\allowbreak{}c}$\to$\texttt{155\_\allowbreak{}a}~(0.824), \texttt{-10\_\allowbreak{}c}$\to$\texttt{155\_\allowbreak{}a}~(-0.509), \texttt{-16\_\allowbreak{}a}$\to$\texttt{155\_\allowbreak{}a}~(1.11).
\item \textbf{Covers} (coverage, own and as part of the hierarchy): 3 candidate nodes.
\end{itemize}

\textbf{Hypothesis.} Node 155\_\allowbreak{}a computes a 'benign profile' gate: 155\_\allowbreak{}a = ReLU(0.824*sex:0 + 1.109*exang:0 - 0.509*fbs:0 + 0.088). Activates for female patients (sex:0=1) with no exertional angina (exang:0=1); normal fasting glucose slightly attenuates. Output-facing copy of hidden node 155. Activation profile: active 70\% / hard-zero 30\% (verified). Reads out to output with weight -0.960, lowering predicted risk. Ablation: removing this readout drops test AUC 0.897-$>$0.869 (-0.028) - the single largest nonlinear contribution to generalisation.

\textbf{Closed form:}
{\small\begingroup\allowdisplaybreaks
\begin{align*}
y_{0} &= \operatorname{relu}\bigl(0.0879 + 1.11 x_{exang:0\_a} + 0.824 x_{sex:0\_b\_c} \\
&\qquad - 0.509 x_{fbs:0\_b\_c}\bigr)
\end{align*}\endgroup}

where $y_{0}$ is node \texttt{155\_\allowbreak{}a}; $z_{i}$, where shown, is the pre-activation of a sigmoid output $y_{i}$; weights and biases are shown to 3 significant figures.
 On the training split (n = 242), the displayed closed form differs from the forward pass by at most 0.00159, and the exact function by at most 0.

\textbf{Evidence} (3 items; full payloads in the JSON record):
\begin{itemize}
\item \textsc{analytical}, \texttt{analytical\_\allowbreak{}formula}; data split: n/a (analytical). Closed-form gate. A shallow, single-step gate encoding a recognised protective factor (female + no exertional angina).
\item \textsc{empirical}, \texttt{activation\_\allowbreak{}profile}; data split: test (inferred, inferred from narrative). Activation profile verified on the test split: active for \textasciitilde{}70\% of patients, hard-zero for \textasciitilde{}30\%. A genuinely used gate, not a dead or saturated unit.
\item \textsc{empirical}, \texttt{ablation\_\allowbreak{}study}; data split: test (inferred, inferred from payload keys). The single largest nonlinear contribution to generalisation: removing this gate costs 0.028 test AUC. A shallow protective-factor gate that genuinely generalises. Re-verified on the full test split.
\end{itemize}

\textbf{Authorship:} drafted by unknown on 2026-07-08; verified by nobody (unverified).

\subsubsection{\texorpdfstring{A3: restbp\_\allowbreak{}chol\_\allowbreak{}interaction}{A3: restbp chol interaction}}\label{app:heart-paper:ann-restbp-chol-interaction}
% annotation id ann_37, operation seq 37
\textbf{Kind:} Leaf annotation (depth 1); parent A1 (\cref{app:heart-paper:ann-heart-disease-risk-score}).

\begin{itemize}
\item \textbf{Entry nodes:} \texttt{-8} (trestbps), \texttt{-9\_\allowbreak{}c} (chol\_\allowbreak{}c).
\item \textbf{Exit nodes:} \texttt{2379}.
\item \textbf{Internal nodes:} \texttt{3240}.
\item \textbf{Edges (weight):} \texttt{-8}$\to$\texttt{3240}~(0.978), \texttt{3240}$\to$\texttt{2379}~(0.931), \texttt{-9\_\allowbreak{}c}$\to$\texttt{2379}~(-0.529).
\item \textbf{Covers} (coverage, own and as part of the hierarchy): 3 candidate nodes.
\end{itemize}

\textbf{Hypothesis.} Subgraph \{trestbps, chol\_\allowbreak{}c, 3240, 2379\}: a near-linear resting-blood-pressure signal. 3240 = ReLU(0.978*trestbps + 3.4) is always-on (100\% active, verified) - a linear pass-through, not a gate. 2379 = ReLU(0.931*3240 - 0.529*chol\_\allowbreak{}c - 0.754) simplifies to ReLU(0.911*trestbps - 0.529*chol\_\allowbreak{}c + 2.41), near-always-on. Reads out to output with +0.364 (higher resting BP raises risk). Ablation: removing this readout drops test AUC 0.897-$>$0.885 (-0.012) - a modest but real generalising contribution.

\textbf{Closed form:}
{\small\begingroup\allowdisplaybreaks
\begin{align*}
y_{0} &= \operatorname{relu}\bigl(-0.754 + 0.931 \operatorname{relu}\bigl(3.4 \\
&\qquad + 0.978 x_{trestbps}\bigr) - 0.529 x_{chol\_c}\bigr)
\end{align*}\endgroup}

where $y_{0}$ is node \texttt{2379}; $z_{i}$, where shown, is the pre-activation of a sigmoid output $y_{i}$; weights and biases are shown to 3 significant figures.
 On the training split (n = 242), the displayed closed form differs from the forward pass by at most 0.00246, and the exact function by at most $8.88\times 10^{-16}$.

\textbf{Evidence} (3 items; full payloads in the JSON record):
\begin{itemize}
\item \textsc{analytical}, \texttt{analytical\_\allowbreak{}formula}; data split: n/a (analytical). Despite the two-node structure this collapses to a near-linear resting-BP term: 3240 is always-on (a pass-through, not a gate) and 2379 is near-always-on. A shallow, generalising risk term.
\item \textsc{empirical}, \texttt{activation\_\allowbreak{}profile}; data split: not recorded. Activation profile verified: 3240 is 100\% active - confirming it is a linear pass-through rather than a gate, which is why the whole pathway reduces to a near-linear resting-BP term.
\item \textsc{empirical}, \texttt{ablation\_\allowbreak{}study}; data split: test (inferred, inferred from payload keys). A modest but real generalising contribution: removing the resting-BP term costs 0.012 test AUC. Re-verified on the test split.
\end{itemize}

\textbf{Authorship:} drafted by unknown on 2026-07-08; verified by nobody (unverified).

\subsubsection{\texorpdfstring{A4: risk\_\allowbreak{}aggregation\_\allowbreak{}gate}{A4: risk aggregation gate}}\label{app:heart-paper:ann-risk-aggregation-gate}
% annotation id ann_38, operation seq 38
\textbf{Kind:} Composition annotation (depth 1); parent A1 (\cref{app:heart-paper:ann-heart-disease-risk-score}); children A5 (\cref{app:heart-paper:ann-benign-profile-gate-534path}), A6 (\cref{app:heart-paper:ann-chol-threshold}).

\begin{itemize}
\item \textbf{Entry nodes:} \texttt{-22\_\allowbreak{}b} (ca\_\allowbreak{}b), \texttt{-9\_\allowbreak{}b} (chol\_\allowbreak{}b).
\item \textbf{Exit nodes:} \texttt{534}.
\item \textbf{Internal nodes:} none.
\item \textbf{Edges (weight):} \texttt{155\_\allowbreak{}b}$\to$\texttt{534}~(0.818), \texttt{3040}$\to$\texttt{534}~(-0.677), \texttt{-22\_\allowbreak{}b}$\to$\texttt{534}~(-1.16), \texttt{-9\_\allowbreak{}b}$\to$\texttt{534}~(-0.592).
\item \textbf{With its descendants:} 9 nodes.
\item \textbf{Covers} (coverage, own and as part of the hierarchy): 8 candidate nodes.
\end{itemize}

\textbf{Hypothesis.} Node 534 = ReLU(0.818*\allowbreak{}[benign\_\allowbreak{}profile\_\allowbreak{}gate\_\allowbreak{}534path]\allowbreak{} - 0.677*\allowbreak{}[chol\_\allowbreak{}threshold]\allowbreak{} - 1.155*ca - 0.592*chol - 0.761): the model's most elaborate unit, a genuine ReLU gate (hard-zero 53.5\% verified) composing four risk signals - vessel count (ca), a cholesterol copy, the cholesterol threshold (3040), and the redundant benign-profile copy (155\_\allowbreak{}b). Net per-unit output effects: ca +1.02, chol +0.52, chol\_\allowbreak{}threshold +0.60, benign -0.72; when the benign signal is strong the whole block is clipped off together. KEY ABLATION RESULT: disabling 534's readout changes test AUC only 0.897-$>$0.899 (+0.002) while train drops 0.946-$>$0.942 - i.e. this showpiece unit is essentially train-only and does not generalise. Signs match cardiology but the structure is fit to the 242 training cases.

\textbf{Closed form} (children as named functions):
{\small\begingroup\allowdisplaybreaks
\begin{align*}
y_{0} &= \operatorname{relu}\bigl(-0.761 \\
&\qquad + 0.818 \operatorname{benign\_profile\_gate\_534path}\bigl(x_{sex:0\_b\_d}, x_{fbs:0\_b\_d}, \\
&\qquad x_{exang:0\_b}\bigr) - 0.592 x_{chol\_b} - 1.16 x_{ca\_b} \\
&\qquad - 0.677 \operatorname{chol\_threshold}\bigl(x_{chol\_d}\bigr)\bigr)
\end{align*}\endgroup}

Expanded over the children:
{\small\begingroup\allowdisplaybreaks
\begin{align*}
y_{0} &= \operatorname{relu}\bigl(-0.761 + 0.818 \operatorname{relu}\bigl(0.0879 \\
&\qquad + 1.11 x_{exang:0\_b} + 0.824 x_{sex:0\_b\_d} - 0.509 x_{fbs:0\_b\_d}\bigr) \\
&\qquad - 0.592 x_{chol\_b} - 1.16 x_{ca\_b} - 0.677 \operatorname{relu}\bigl(-0.864 \\
&\qquad + 1.05 x_{chol\_d}\bigr)\bigr)
\end{align*}\endgroup}

where $y_{0}$ is node \texttt{534}; $z_{i}$, where shown, is the pre-activation of a sigmoid output $y_{i}$; weights and biases are shown to 3 significant figures.
 On the training split (n = 242), the displayed closed form differs from the forward pass by at most 0.00694, and the exact function by at most $4.44\times 10^{-16}$.

\textbf{Evidence} (3 items; full payloads in the JSON record):
\begin{itemize}
\item \textsc{analytical}, \texttt{analytical\_\allowbreak{}formula}; data split: n/a (analytical). Composed closed form (the composition annotation resolves into its two children plus two direct inputs). On paper the model's most sophisticated structure - a rectified gate folding four risk signals together. Signs match cardiology.
\item \textsc{empirical}, \texttt{activation\_\allowbreak{}profile}; data split: not recorded. Activation profile verified: hard-zero for 53.5\% of patients, confirming 534 is a genuine gate rather than a pass-through.
\item \textsc{empirical}, \texttt{ablation\_\allowbreak{}study}; data split: train+test (inferred, inferred from payload keys). KEY RESULT. The model's most elaborate structure does no real work on held-out data: removing it changes test AUC by +0.002 (slightly helps) while costing on train. Structural elaboration is where the overfitting concentrates. Only holding structure against ablation reveals it - the gate looks like the centrepiece, yet measurement shows it is decorative while the plainest gates do the work. Re-verified on both splits.
\end{itemize}

\textbf{Authorship:} drafted by unknown on 2026-07-08; verified by nobody (unverified).

\subsubsection{\texorpdfstring{A5: benign\_\allowbreak{}profile\_\allowbreak{}gate\_\allowbreak{}534path}{A5: benign profile gate 534path}}\label{app:heart-paper:ann-benign-profile-gate-534path}
% annotation id ann_35, operation seq 35
\textbf{Kind:} Leaf annotation (depth 2); parent A4 (\cref{app:heart-paper:ann-risk-aggregation-gate}).

\begin{itemize}
\item \textbf{Entry nodes:} \texttt{-2\_\allowbreak{}d} (sex:0\_\allowbreak{}b\_\allowbreak{}d), \texttt{-10\_\allowbreak{}d} (fbs:0\_\allowbreak{}b\_\allowbreak{}d), \texttt{-16\_\allowbreak{}b} (exang:0\_\allowbreak{}b).
\item \textbf{Exit nodes:} \texttt{155\_\allowbreak{}b}.
\item \textbf{Internal nodes:} none.
\item \textbf{Edges (weight):} \texttt{-2\_\allowbreak{}d}$\to$\texttt{155\_\allowbreak{}b}~(0.824), \texttt{-10\_\allowbreak{}d}$\to$\texttt{155\_\allowbreak{}b}~(-0.509), \texttt{-16\_\allowbreak{}b}$\to$\texttt{155\_\allowbreak{}b}~(1.11).
\item \textbf{Covers} (coverage, own and as part of the hierarchy): 3 candidate nodes.
\end{itemize}

\textbf{Hypothesis.} Node 155\_\allowbreak{}b computes the SAME benign-profile gate as 155\_\allowbreak{}a (identical incoming weights and bias 0.088); it is the second copy of hidden node 155, routed into node 534 instead of the output. The benign signal is therefore encoded twice - once directly to the output (155\_\allowbreak{}a) and once inside the 534 gate (155\_\allowbreak{}b). 155\_\allowbreak{}b enters 534 with weight +0.818; net path effect (+0.818)*(-0.884) = -0.72 per unit, same risk-lowering direction as the direct copy.

\textbf{Closed form:}
{\small\begingroup\allowdisplaybreaks
\begin{align*}
y_{0} &= \operatorname{relu}\bigl(0.0879 + 1.11 x_{exang:0\_b} + 0.824 x_{sex:0\_b\_d} \\
&\qquad - 0.509 x_{fbs:0\_b\_d}\bigr)
\end{align*}\endgroup}

where $y_{0}$ is node \texttt{155\_\allowbreak{}b}; $z_{i}$, where shown, is the pre-activation of a sigmoid output $y_{i}$; weights and biases are shown to 3 significant figures.
 On the training split (n = 242), the displayed closed form differs from the forward pass by at most 0.00159, and the exact function by at most 0.

\textbf{Evidence} (1 item; full payloads in the JSON record):
\begin{itemize}
\item \textsc{analytical}, \texttt{analytical\_\allowbreak{}formula}; data split: n/a (analytical). 155\_\allowbreak{}b shares incoming weights and bias 0.088 with 155\_\allowbreak{}a exactly (bit-identical). The benign protective signal is therefore redundantly encoded twice - once direct to output, once inside the 534 gate. Redundant double-encoding is a symptom of the same overfitting the 534 gate exhibits.
\end{itemize}

\textbf{Authorship:} drafted by unknown on 2026-07-08; verified by nobody (unverified).

\subsubsection{\texorpdfstring{A6: chol\_\allowbreak{}threshold}{A6: chol threshold}}\label{app:heart-paper:ann-chol-threshold}
% annotation id ann_36, operation seq 36
\textbf{Kind:} Leaf annotation (depth 2); parent A4 (\cref{app:heart-paper:ann-risk-aggregation-gate}).

\begin{itemize}
\item \textbf{Entry nodes:} \texttt{-9\_\allowbreak{}d} (chol\_\allowbreak{}d).
\item \textbf{Exit nodes:} \texttt{3040}.
\item \textbf{Internal nodes:} none.
\item \textbf{Edges (weight):} \texttt{-9\_\allowbreak{}d}$\to$\texttt{3040}~(1.05).
\item \textbf{Covers} (coverage, own and as part of the hierarchy): 1 candidate node.
\end{itemize}

\textbf{Hypothesis.} Node 3040 = ReLU(1.052*chol - 0.864): a one-sided cholesterol threshold. Since cholesterol is z-normalised, this turns on only above \textasciitilde{}0.82 SD above the mean. Activation profile: hard-zero 82.5\% / active 17.5\% (verified) - a sparse threshold detector. Feeds node 534 with -0.677; net path effect +0.60 per unit (high cholesterol raises predicted risk). This is the sub-unit that lets cholesterol act as a risk factor only above a level rather than linearly - but it sits inside the 534 gate, which the ablation shows contributes nothing on test.

\textbf{Closed form:}
{\small\begingroup\allowdisplaybreaks
\begin{align*}
y_{0} &= \operatorname{relu}\bigl(-0.864 + 1.05 x_{chol\_d}\bigr)
\end{align*}\endgroup}

where $y_{0}$ is node \texttt{3040}; $z_{i}$, where shown, is the pre-activation of a sigmoid output $y_{i}$; weights and biases are shown to 3 significant figures.
 On the training split (n = 242), the displayed closed form differs from the forward pass by at most 0.00824, and the exact function by at most 0.

\textbf{Evidence} (2 items; full payloads in the JSON record):
\begin{itemize}
\item \textsc{analytical}, \texttt{analytical\_\allowbreak{}formula}; data split: test (inferred, inferred from narrative). A sparse one-sided cholesterol threshold - the sub-unit that lets cholesterol act as a risk factor only above a level rather than linearly. But it sits inside the 534 gate, which ablation shows contributes nothing on test.
\item \textsc{empirical}, \texttt{activation\_\allowbreak{}profile}; data split: not recorded. Activation profile verified: fires for only \textasciitilde{}17.5\% of patients, confirming a sparse one-sided threshold detector rather than a linear response.
\end{itemize}

\textbf{Authorship:} drafted by unknown on 2026-07-08; verified by nobody (unverified).

\subsection{Audit against the definitions}\label{app:heart-paper:audit}
Automatic findings on the stored explanation (codes as in the JSON record's \texttt{issues}); they flag what an examiner should check, not verdicts.
\begin{itemize}
\item \textsc{warning} \texttt{unverified}: No human verification is recorded for this explanation.
\item \textsc{warning} \texttt{evidence\_\allowbreak{}without\_\allowbreak{}split} [chol\_\allowbreak{}threshold]: Empirical evidence item 1 (activation\_\allowbreak{}profile) records no data split.
\item \textsc{warning} \texttt{evidence\_\allowbreak{}without\_\allowbreak{}split} [restbp\_\allowbreak{}chol\_\allowbreak{}interaction]: Empirical evidence item 1 (activation\_\allowbreak{}profile) records no data split.
\item \textsc{warning} \texttt{evidence\_\allowbreak{}without\_\allowbreak{}split} [risk\_\allowbreak{}aggregation\_\allowbreak{}gate]: Empirical evidence item 1 (activation\_\allowbreak{}profile) records no data split.
\item \textsc{info} \texttt{authorship\_\allowbreak{}unknown}: Drafter of the explanation is not recorded.
\item \textsc{info} \texttt{evidence\_\allowbreak{}split\_\allowbreak{}inferred} (6 findings, annotations benign\_\allowbreak{}profile\_\allowbreak{}gate, heart\_\allowbreak{}disease\_\allowbreak{}risk\_\allowbreak{}score, restbp\_\allowbreak{}chol\_\allowbreak{}interaction, risk\_\allowbreak{}aggregation\_\allowbreak{}gate).
\end{itemize}
\end{drafted}
